\documentclass[10pt,a4paper]{article}

\usepackage[utf8]{inputenc}
\usepackage[T1]{fontenc}
\usepackage[a4paper,margin=0.6in,top=0.52in,bottom=0.52in]{geometry}
\usepackage{amsmath}
\usepackage{tikz}
\usetikzlibrary{shapes,arrows,positioning,calc,fit,backgrounds}
\usepackage{xcolor}
\usepackage{lmodern}
\usepackage{booktabs}
\usepackage{tabularx}
\usepackage{enumitem}
\usepackage{titlesec}
\usepackage{fancyhdr}
\usepackage{listings}
\usepackage{microtype}
\usepackage[colorlinks=true,linkcolor=primary,urlcolor=secondary]{hyperref}

% =========================================================
% Palette & Styling (Matching lab reports)
% =========================================================
\definecolor{primary}{RGB}{24, 43, 73}        % Deep clean navy
\definecolor{secondary}{RGB}{14, 116, 144}   % Teal accent
\definecolor{darkslate}{RGB}{30, 41, 59}     % Normal text color
\definecolor{boxbg}{RGB}{248, 250, 252}      % Soft light box
\definecolor{bordergray}{RGB}{226, 232, 240} % Light border

% Terminal Color Scheme: Modern Light Grey Theme
\definecolor{termbg}{RGB}{246, 248, 250}     % Light grey terminal body
\definecolor{termbar}{RGB}{229, 231, 235}    % Medium light grey title bar
\definecolor{termborder}{RGB}{209, 213, 219} % Crisp border line
\definecolor{termtext}{RGB}{31, 41, 55}      % Dark charcoal text
\definecolor{termuser}{RGB}{22, 101, 52}     % Dark green for user
\definecolor{termhost}{RGB}{21, 128, 61}     % Green for host
\definecolor{termpath}{RGB}{29, 78, 216}     % Blue for directory path
\definecolor{termcmd}{RGB}{17, 24, 39}       % Bold black for commands
\definecolor{termout}{RGB}{55, 65, 81}       % Clear dark grey for outputs
\definecolor{termdir}{RGB}{37, 99, 235}      % Blue for directories
\definecolor{termfile}{RGB}{31, 41, 55}      % Dark for regular files
\definecolor{termhidden}{RGB}{107, 114, 128} % Muted grey
\definecolor{termsize}{RGB}{180, 83, 9}      % Amber for sizes
\definecolor{termperm}{RGB}{79, 70, 229}     % Indigo for permissions

% Code Listing Colors
\definecolor{codebg}{RGB}{248, 250, 252}
\definecolor{codeborder}{RGB}{218, 224, 233}
\definecolor{codekw}{RGB}{14, 116, 144}
\definecolor{codestr}{RGB}{22, 101, 52}
\definecolor{codecomment}{RGB}{100, 116, 139}
\definecolor{codenum}{RGB}{148, 163, 184}

% Table helpers
\newcolumntype{Y}{>{\centering\arraybackslash}X}

% =========================================================
% Section Headings
% =========================================================
\titleformat{\section}
  {\color{primary}\large\bfseries\sffamily\raggedright}
  {\thesection.}{0.5em}{}
\titlespacing*{\section}{0pt}{7pt}{2pt}

\titleformat{\subsection}
  {\color{secondary}\normalsize\bfseries\sffamily\raggedright}
  {\thesubsection}{0.5em}{}
\titlespacing*{\subsection}{0pt}{5pt}{2pt}

% =========================================================
% Headers and Footers
% =========================================================
\pagestyle{fancy}
\fancyhf{}
\renewcommand{\headrulewidth}{0.4pt}
\renewcommand{\footrulewidth}{0pt}
\fancyhead[L]{\small\sffamily\color{darkslate} Operating Systems Lab}
\fancyhead[R]{\small\sffamily\color{darkslate} Appendix: Source Code \& Execution Output}
\fancyfoot[C]{\small\sffamily\thepage}

\fancypagestyle{firstpage}{%
  \fancyhf{}
  \renewcommand{\headrulewidth}{0pt}
  \renewcommand{\footrulewidth}{0pt}
  \fancyfoot[C]{\small\sffamily\thepage}
}

% =========================================================
% Terminal Window Macro
% =========================================================
\newcommand{\termprompt}[1][\textasciitilde]{%
  \textbf{\textcolor{termuser}{ratul}}\textcolor{termtext}{@}\textbf{\textcolor{termhost}{firefly}}\textcolor{termtext}{:}\textbf{\textcolor{termpath}{#1}}\textcolor{termcmd}{\$}\ %
}

\newcommand{\terminalwindow}[3][\dimexpr\linewidth-0.6pt\relax]{%
\par\vspace{2pt}%
{\centering
\begin{tikzpicture}
  % 1. Terminal text on main layer
  \node[
      text=termtext,
      font=\ttfamily\fontsize{7.2}{9.2}\selectfont,
      align=left,
      text width=\dimexpr#1-1.02cm\relax,
      inner xsep=6pt,
      inner ysep=6pt,
      anchor=north west
  ] (body) at (0.32cm, -0.56cm) {#3};

  % Coordinates
  \coordinate (topleft) at (0, 0);
  \coordinate (bottomright) at ($(body.south -| #1, 0) + (0, -5pt)$);
  \coordinate (titlecenter) at ($0.5*(topleft) + 0.5*(#1, 0) + (0, -0.28cm)$);

  % 2. Background and Title bar on background layer
  \begin{scope}[on background layer]
    \fill[termbg, rounded corners=6pt] (topleft) rectangle (bottomright);

    \begin{scope}
      \clip[rounded corners=6pt] (topleft) rectangle (bottomright);
      \fill[termbar] (topleft) rectangle (#1, -0.56cm);
      \draw[termborder, line width=0.5pt] (0, -0.56cm) -- (#1, -0.56cm);
    \end{scope}

    \draw[termborder, line width=0.6pt, rounded corners=6pt] (topleft) rectangle (bottomright);
  \end{scope}

  % 3. Window control dots
  \fill[red!80!black] (0.36cm, -0.28cm) circle (0.07cm);
  \fill[yellow!80!black] (0.63cm, -0.28cm) circle (0.07cm);
  \fill[green!75!black] (0.90cm, -0.28cm) circle (0.07cm);

  % 4. Centered terminal title
  \node[
      text=gray!40!black,
      font=\sffamily\bfseries\fontsize{7.2}{8}\selectfont,
      anchor=center
  ] at (titlecenter) {#2};
\end{tikzpicture}\par}%
\vspace{2pt}%
}

\newcommand{\cmd}[1]{\texttt{\textbf{#1}}}

% =========================================================
% Code Listing Configuration
% =========================================================
\lstset{
  backgroundcolor=\color{codebg},
  basicstyle=\ttfamily\fontsize{6.9}{8.4}\selectfont\color{darkslate},
  keywordstyle=\bfseries\color{codekw},
  commentstyle=\itshape\color{codecomment},
  stringstyle=\color{codestr},
  numbers=left,
  numberstyle=\tiny\color{codenum},
  stepnumber=1,
  numbersep=8pt,
  showstringspaces=false,
  breaklines=true,
  frame=single,
  rulecolor=\color{codeborder},
  tabsize=4,
  xleftmargin=14pt,
  framexleftmargin=14pt
}

\begin{document}
\raggedbottom
\thispagestyle{firstpage}

% =========================================================
% Title Block
% =========================================================
\begin{center}
  {\Large\bfseries\sffamily\color{primary} Operating Systems Lab: Source Code and Program Outputs}\\[3pt]
  {\normalsize\sffamily\textbf{Attachment / Appendix:} CPU Scheduling Implementations and Executions (Labs 8, 9, and 10)}
\end{center}
\vspace{-4pt}

% ==============================================================================
% LAB 8
% ==============================================================================
\section{Lab 8: First Come, First Served (FCFS) Scheduling}

\subsection{Source Code (C++)}
\begin{lstlisting}[language=C++]
#include <bits/stdc++.h>
using namespace std;

struct Process {
    int id;
    int arrival;
    int execution;
    int start;
    int completion;
    int turnaround;
    int waiting;
};

int main() {
    int n;

    cout << "Enter the number of processes: ";
    cin >> n;

    if (cin.fail() || n <= 0) {
        cout << "Invalid number of processes!\n";
        return 1;
    }

    vector<Process> processes(n);

    // Input process details
    for (int i = 0; i < n; i++) {
        processes[i].id = i + 1;

        cout << "Enter arrival and execution time for process "
             << processes[i].id << ": ";

        cin >> processes[i].arrival >> processes[i].execution;

        if (cin.fail() ||
            processes[i].arrival < 0 ||
            processes[i].execution <= 0) {
            cout << "Invalid input! Arrival time must be "
                 << "non-negative and execution time must "
                 << "be positive.\n";
            return 1;
        }
    }

    // Sort by arrival time; preserve input order for ties
    stable_sort(processes.begin(), processes.end(),
        [](const Process& a, const Process& b) {
            return a.arrival < b.arrival;
        });

    int currentTime = 0;
    double totalWaiting = 0;
    double totalTurnaround = 0;

    // Calculate scheduling times
    for (Process& p : processes) {
        // Handle CPU idle time
        currentTime = max(currentTime, p.arrival);

        p.start = currentTime;
        p.waiting = p.start - p.arrival;

        currentTime += p.execution;

        p.completion = currentTime;
        p.turnaround = p.completion - p.arrival;

        totalWaiting += p.waiting;
        totalTurnaround += p.turnaround;
    }

    // Display results in original process order
    sort(processes.begin(), processes.end(),
        [](const Process& a, const Process& b) {
            return a.id < b.id;
        });

    cout << "\nProcess\tArrival\tExecution\tStart\tCompletion"
         << "\tTurnaround\tWaiting\n";

    for (const Process& p : processes) {
        cout << "P" << p.id << "\t"
             << p.arrival << "\t"
             << p.execution << "\t\t"
             << p.start << "\t"
             << p.completion << "\t"
             << p.turnaround << "\t\t"
             << p.waiting << '\n';
    }

    // Display average times
    cout << fixed << setprecision(2);

    cout << "\nAverage Waiting Time: "
         << totalWaiting / n << '\n';

    cout << "Average Turnaround Time: "
         << totalTurnaround / n << '\n';

    // Display Gantt chart in FCFS execution order
    stable_sort(processes.begin(), processes.end(),
        [](const Process& a, const Process& b) {
            return a.arrival < b.arrival;
        });

    cout << "\nGantt Chart:\n";

    for (const Process& p : processes) {
        if (p.start > 0 && p.start > p.arrival) {
            // The CPU may have been idle before this process.
            // Idle intervals are displayed below using timestamps.
        }

        cout << "| ";
        if (p.start > 0) {
            cout << "P" << p.id << " ";
        } else {
            cout << "P" << p.id << " ";
        }
    }
    cout << "|\n";

    // Print process start and completion timestamps
    if (!processes.empty()) {
        cout << processes[0].start;

        for (const Process& p : processes) {
            cout << "\t" << p.completion;
        }

        cout << '\n';
    }

    return 0;
}
\end{lstlisting}

\subsection{Execution Output}
\terminalwindow{ratul@firefly: \textasciitilde/OS-Lab\ \textemdash\ bash}{
% Command 1: compile
\termprompt[\textasciitilde/OS-Lab]\textbf{\textcolor{termcmd}{g++ fcfs.cpp -o fcfs}}

\vspace{0.06cm}
% Command 2: run
\termprompt[\textasciitilde/OS-Lab]\textbf{\textcolor{termcmd}{./fcfs}}

\textcolor{termout}{Enter the number of processes: 4}

\textcolor{termout}{Enter arrival and execution time for process 1: 0 4}

\textcolor{termout}{Enter arrival and execution time for process 2: 1 3}

\textcolor{termout}{Enter arrival and execution time for process 3: 2 1}

\textcolor{termout}{Enter arrival and execution time for process 4: 3 2}

\vspace{0.06cm}
\textcolor{termperm}{\textbf{Process\ \ Arrival\ \ Execution\ \ Start\ \ Completion\ \ Turnaround\ \ Waiting}}

\textcolor{termout}{P1\ \ \ \ \ \ \ 0\ \ \ \ \ \ \ \ 4\ \ \ \ \ \ \ \ \ \ 0\ \ \ \ \ \ 4\ \ \ \ \ \ \ \ \ \ \ 4\ \ \ \ \ \ \ \ \ \ \ 0}

\textcolor{termout}{P2\ \ \ \ \ \ \ 1\ \ \ \ \ \ \ \ 3\ \ \ \ \ \ \ \ \ \ 4\ \ \ \ \ \ 7\ \ \ \ \ \ \ \ \ \ \ 6\ \ \ \ \ \ \ \ \ \ \ 3}

\textcolor{termout}{P3\ \ \ \ \ \ \ 2\ \ \ \ \ \ \ \ 1\ \ \ \ \ \ \ \ \ \ 7\ \ \ \ \ \ 8\ \ \ \ \ \ \ \ \ \ \ 6\ \ \ \ \ \ \ \ \ \ \ 5}

\textcolor{termout}{P4\ \ \ \ \ \ \ 3\ \ \ \ \ \ \ \ 2\ \ \ \ \ \ \ \ \ \ 8\ \ \ \ \ \ 10\ \ \ \ \ \ \ \ \ \ 7\ \ \ \ \ \ \ \ \ \ \ 5}

\vspace{0.06cm}
\textcolor{termuser}{\textbf{Average Waiting Time: 3.25}}

\textcolor{termuser}{\textbf{Average Turnaround Time: 5.75}}

\vspace{0.06cm}
\textcolor{termsize}{\textbf{Gantt Chart:}}

\textcolor{termcmd}{| P1 | P2 | P3 | P4 |}

\textcolor{termout}{0\ \ \ \ \ \ \ 4\ \ \ \ \ \ \ 7\ \ \ \ \ \ \ 8\ \ \ \ \ \ \ 10}

\vspace{0.08cm}
% Final prompt
\termprompt[\textasciitilde/OS-Lab]\textbf{\textcolor{termcmd}{\rule{0pt}{0.9em}\textunderscore}}
}

\newpage

% ==============================================================================
% LAB 9
% ==============================================================================
\section{Lab 9: Shortest Job First (SJF Non-Preemptive) Scheduling}

\subsection{Source Code (C)}
\begin{lstlisting}[language=C]
#include <stdio.h>

int main() {
    int n;

    printf("Enter the number of processes: ");
    scanf("%d", &n);

    int process[n], arrival[n], burst[n];
    int completion[n], turnaround[n], waiting[n];
    int is_completed[n];

    int i, j, temp, pos;
    int current_time = 0;
    int completed = 0;

    float totalTAT = 0;
    float totalWT = 0;
    float averageTAT, averageWT;

    // Input process details
    for (i = 0; i < n; i++) {
        process[i] = i + 1;

        printf("Enter arrival time of process %d: ", process[i]);
        scanf("%d", &arrival[i]);

        printf("Enter burst time of process %d: ", process[i]);
        scanf("%d", &burst[i]);

        completion[i] = -1;
        is_completed[i] = 0;
    }

    // Sorting processes based on arrival time
    for (i = 0; i < n; i++) {
        pos = i;

        for (j = i + 1; j < n; j++) {
            if (arrival[j] < arrival[pos]) {
                pos = j;
            }
        }

        // Swap arrival time
        temp = arrival[i];
        arrival[i] = arrival[pos];
        arrival[pos] = temp;

        // Swap burst time
        temp = burst[i];
        burst[i] = burst[pos];
        burst[pos] = temp;

        // Swap process ID
        temp = process[i];
        process[i] = process[pos];
        process[pos] = temp;
    }

    // Apply Shortest Job First (Non-Preemptive)
    while (completed != n) {

        int min_burst = 9999;
        int index = -1;

        // Find the shortest available process
        for (i = 0; i < n; i++) {

            if (arrival[i] <= current_time && is_completed[i] == 0) {

                if (burst[i] < min_burst) {
                    min_burst = burst[i];
                    index = i;
                }

                // If burst times are equal, select the one
                // that arrived first
                else if (burst[i] == min_burst &&
                         arrival[i] < arrival[index]) {
                    index = i;
                }
            }
        }

        // If a process is available
        if (index != -1) {

            completion[index] = current_time + burst[index];
            turnaround[index] = completion[index] - arrival[index];
            waiting[index] = turnaround[index] - burst[index];

            // Add to totals
            totalTAT = totalTAT + turnaround[index];
            totalWT = totalWT + waiting[index];

            // Update current time
            current_time = completion[index];

            is_completed[index] = 1;
            completed++;
        }

        // If no process has arrived yet
        else {
            current_time++;
        }
    }

    // Calculate averages
    averageTAT = totalTAT / n;
    averageWT = totalWT / n;

    // Printing the final process table
    printf("\nProcess\tArrival Time\tBurst Time\tCompletion Time\tTurnaround Time\tWaiting Time\n");

    for (i = 0; i < n; i++) {
        printf("P%d\t%d\t\t%d\t\t%d\t\t%d\t\t%d\n",
               process[i],
               arrival[i],
               burst[i],
               completion[i],
               turnaround[i],
               waiting[i]);
    }

    // Print total and average values
    printf("\nTotal Turnaround Time = %.0f", totalTAT);
    printf("\nTotal Waiting Time    = %.0f", totalWT);

    printf("\nAverage Turnaround Time = %.2f", averageTAT);
    printf("\nAverage Waiting Time    = %.2f\n", averageWT);

    return 0;
}
\end{lstlisting}

\subsection{Execution Output}
\terminalwindow{ratul@firefly: \textasciitilde/OS-Lab\ \textemdash\ bash}{
% Command 1: compile
\termprompt[\textasciitilde/OS-Lab]\textbf{\textcolor{termcmd}{gcc sjf.c -o sjf}}

\vspace{0.06cm}
% Command 2: run
\termprompt[\textasciitilde/OS-Lab]\textbf{\textcolor{termcmd}{./sjf}}

\textcolor{termout}{Enter the number of processes: 4}

\textcolor{termout}{Enter arrival time of process 1: 0}

\textcolor{termout}{Enter burst time of process 1: 7}

\textcolor{termout}{Enter arrival time of process 2: 2}

\textcolor{termout}{Enter burst time of process 2: 4}

\textcolor{termout}{Enter arrival time of process 3: 4}

\textcolor{termout}{Enter burst time of process 3: 1}

\textcolor{termout}{Enter arrival time of process 4: 5}

\textcolor{termout}{Enter burst time of process 4: 4}

\vspace{0.06cm}
\textcolor{termperm}{\textbf{Process\ \ Arrival\ Time\ \ Burst\ Time\ \ Completion\ Time\ \ Turnaround\ Time\ \ Waiting\ Time}}

\textcolor{termout}{P1\ \ \ \ \ \ \ 0\ \ \ \ \ \ \ \ \ \ \ \ \ 7\ \ \ \ \ \ \ \ \ \ \ 7\ \ \ \ \ \ \ \ \ \ \ \ \ \ \ \ 7\ \ \ \ \ \ \ \ \ \ \ \ \ \ \ \ 0}

\textcolor{termout}{P2\ \ \ \ \ \ \ 2\ \ \ \ \ \ \ \ \ \ \ \ \ 4\ \ \ \ \ \ \ \ \ \ \ 12\ \ \ \ \ \ \ \ \ \ \ \ \ \ 10\ \ \ \ \ \ \ \ \ \ \ \ \ \ \ 6}

\textcolor{termout}{P3\ \ \ \ \ \ \ 4\ \ \ \ \ \ \ \ \ \ \ \ \ 1\ \ \ \ \ \ \ \ \ \ \ 8\ \ \ \ \ \ \ \ \ \ \ \ \ \ \ \ 4\ \ \ \ \ \ \ \ \ \ \ \ \ \ \ \ 3}

\textcolor{termout}{P4\ \ \ \ \ \ \ 5\ \ \ \ \ \ \ \ \ \ \ \ \ 4\ \ \ \ \ \ \ \ \ \ \ 16\ \ \ \ \ \ \ \ \ \ \ \ \ \ 11\ \ \ \ \ \ \ \ \ \ \ \ \ \ \ 7}

\vspace{0.06cm}
\textcolor{termsize}{\textbf{Total Turnaround Time = 32}}

\textcolor{termsize}{\textbf{Total Waiting Time\ \ \ \ = 16}}

\vspace{0.04cm}
\textcolor{termuser}{\textbf{Average Turnaround Time = 8.00}}

\textcolor{termuser}{\textbf{Average Waiting Time\ \ \ \ = 4.00}}

\vspace{0.08cm}
% Final prompt
\termprompt[\textasciitilde/OS-Lab]\textbf{\textcolor{termcmd}{\rule{0pt}{0.9em}\textunderscore}}
}

\newpage

% ==============================================================================
% LAB 10
% ==============================================================================
\section{Lab 10: Shortest Remaining Job First (SRJF Preemptive) Scheduling}

\subsection{Source Code (C)}
\begin{lstlisting}[language=C]
#include <stdio.h>

int main() {
    int n, i, j, temp, pos, current_time = 0, completed = 0;
    float totalTAT = 0, totalWT = 0, averageTAT, averageWT;

    printf("Enter the number of processes: ");
    scanf("%d", &n);

    int process[n], arrival[n], burst[n], remaining[n], completion[n], turnaround[n], waiting[n], is_completed[n];

    // Input process details
    for (i = 0; i < n; i++) {
        process[i] = i + 1;
        printf("Enter arrival time of process %d: ", process[i]);
        scanf("%d", &arrival[i]);
        printf("Enter burst time of process %d: ", process[i]);
        scanf("%d", &burst[i]);
        remaining[i] = burst[i];
        completion[i] = -1;
        is_completed[i] = 0;
    }

    // Sort processes based on arrival time
    for (i = 0; i < n; i++) {
        pos = i;
        for (j = i + 1; j < n; j++) {
            if (arrival[j] < arrival[pos]) pos = j;
        }
        temp = arrival[i]; arrival[i] = arrival[pos]; arrival[pos] = temp;
        temp = burst[i]; burst[i] = burst[pos]; burst[pos] = temp;
        temp = remaining[i]; remaining[i] = remaining[pos]; remaining[pos] = temp;
        temp = process[i]; process[i] = process[pos]; process[pos] = temp;
    }

    // Apply Shortest Remaining Time First (SRTF / Preemptive SJF)
    while (completed != n) {
        int min_remaining = 9999, index = -1;

        // Find process with shortest remaining time among arrived processes
        for (i = 0; i < n; i++) {
            if (arrival[i] <= current_time && is_completed[i] == 0) {
                if (remaining[i] < min_remaining && remaining[i] > 0) {
                    min_remaining = remaining[i];
                    index = i;
                } else if (remaining[i] == min_remaining && arrival[i] < arrival[index]) {
                    index = i;
                }
            }
        }

        if (index != -1) {
            remaining[index]--;
            current_time++;

            if (remaining[index] == 0) {
                completion[index] = current_time;
                turnaround[index] = completion[index] - arrival[index];
                waiting[index] = turnaround[index] - burst[index];
                totalTAT += turnaround[index];
                totalWT += waiting[index];
                is_completed[index] = 1;
                completed++;
            }
        } else {
            current_time++;
        }
    }

    averageTAT = totalTAT / n;
    averageWT = totalWT / n;

    // Printing final process table
    printf("\nProcess\tArrival Time\tBurst Time\tCompletion Time\tTurnaround Time\tWaiting Time\n");
    for (i = 0; i < n; i++) {
        printf("P%d\t%d\t\t%d\t\t%d\t\t%d\t\t%d\n",
               process[i], arrival[i], burst[i], completion[i], turnaround[i], waiting[i]);
    }

    printf("\nTotal Turnaround Time: %.0f", totalTAT);
    printf("\nTotal Waiting Time: %.0f", totalWT);
    printf("\nAverage Turnaround Time: %.2f", averageTAT);
    printf("\nAverage Waiting Time: %.2f\n", averageWT);

    return 0;
}
\end{lstlisting}

\subsection{Execution Output}
\terminalwindow{ratul@firefly: \textasciitilde/OS-Lab\ \textemdash\ bash}{
% Command 1: compile
\termprompt[\textasciitilde/OS-Lab]\textbf{\textcolor{termcmd}{gcc srtf.c -o srtf}}

\vspace{0.06cm}
% Command 2: run
\termprompt[\textasciitilde/OS-Lab]\textbf{\textcolor{termcmd}{./srtf}}

\textcolor{termout}{Enter the number of processes: 4}

\textcolor{termout}{Enter arrival time of process 1: 0}

\textcolor{termout}{Enter burst time of process 1: 7}

\textcolor{termout}{Enter arrival time of process 2: 2}

\textcolor{termout}{Enter burst time of process 2: 4}

\textcolor{termout}{Enter arrival time of process 3: 4}

\textcolor{termout}{Enter burst time of process 3: 1}

\textcolor{termout}{Enter arrival time of process 4: 5}

\textcolor{termout}{Enter burst time of process 4: 4}

\vspace{0.06cm}
\textcolor{termperm}{\textbf{Process\ \ Arrival\ Time\ \ Burst\ Time\ \ Completion\ Time\ \ Turnaround\ Time\ \ Waiting\ Time}}

\textcolor{termout}{P1\ \ \ \ \ \ \ 0\ \ \ \ \ \ \ \ \ \ \ \ \ 7\ \ \ \ \ \ \ \ \ \ \ 16\ \ \ \ \ \ \ \ \ \ \ \ \ \ 16\ \ \ \ \ \ \ \ \ \ \ \ \ \ 9}

\textcolor{termout}{P2\ \ \ \ \ \ \ 2\ \ \ \ \ \ \ \ \ \ \ \ \ 4\ \ \ \ \ \ \ \ \ \ \ 7\ \ \ \ \ \ \ \ \ \ \ \ \ \ \ 5\ \ \ \ \ \ \ \ \ \ \ \ \ \ \ 1}

\textcolor{termout}{P3\ \ \ \ \ \ \ 4\ \ \ \ \ \ \ \ \ \ \ \ \ 1\ \ \ \ \ \ \ \ \ \ \ 5\ \ \ \ \ \ \ \ \ \ \ \ \ \ \ 1\ \ \ \ \ \ \ \ \ \ \ \ \ \ \ 0}

\textcolor{termout}{P4\ \ \ \ \ \ \ 5\ \ \ \ \ \ \ \ \ \ \ \ \ 4\ \ \ \ \ \ \ \ \ \ \ 11\ \ \ \ \ \ \ \ \ \ \ \ \ \ 6\ \ \ \ \ \ \ \ \ \ \ \ \ \ \ 2}

\vspace{0.06cm}
\textcolor{termsize}{\textbf{Total Turnaround Time: 28}}

\textcolor{termsize}{\textbf{Total Waiting Time\ \ \ \ : 12}}

\vspace{0.04cm}
\textcolor{termuser}{\textbf{Average Turnaround Time = 7.00}}

\textcolor{termuser}{\textbf{Average Waiting Time\ \ \ \ = 3.00}}

\vspace{0.08cm}
% Final prompt
\termprompt[\textasciitilde/OS-Lab]\textbf{\textcolor{termcmd}{\rule{0pt}{0.9em}\textunderscore}}
}

\end{document}
